\documentclass{exam}

\usepackage{amsmath,amssymb,amsfonts}
\usepackage{graphicx}
\usepackage{amsthm}
\usepackage[french]{babel}
\usepackage{enumitem}

\setlength{\topmargin}{-.5in} \setlength{\textheight}{9.25in}
\setlength{\oddsidemargin}{0in} \setlength{\textwidth}{6.8in}

\makeatletter
\def\gap#1{%
  \setbox\@tempboxa\hbox{{\color{white}#1}}%
  \@tempdima\wd\@tempboxa
  \rlap{\unhbox\@tempboxa}%
  \hb@xt@\@tempdima{\dotfill}%
}
\makeatother   

\theoremstyle{definition}
\newtheorem{exercice}{Exercice}

\begin{document}
\Large

\noindent{\bf Contrôle 2 -- 1h15. \hfill Terminales -- Math expertes \hfill }

\medskip\hrule
\vspace{2em}

\begin{exercice}[2 points, 5 min]
Soit $z = 1 - i$, calculez son module et un argument, puis écrivez-le sous sa forme trigonométrique.
\end{exercice}

\begin{exercice}[2 points, 5 min]
On considère \( z = e^{i\frac{\pi}{6}} \). Donnez la forme algébrique de $z$ et $z^{12}$.
\end{exercice}

\begin{exercice}[2 points, 5 min]
Résoudre dans \( \mathbb{C} \) l'équation \( z^{4} = 16 \).
\end{exercice}

% Exercice 23
\begin{exercice}[4 points, 15 min]
\begin{enumerate}
    \item Développer \( (a + b)^4 \).
    \item En utilisant une formule d'Euler, prouver que pour tout réel \( x \),
       \[
       \sin^4(x) = \dfrac{1}{8} \left( \cos(4x) - 4\cos(2x) + 3 \right).
       \]
\end{enumerate}
\end{exercice}

% Exercice 24
\begin{exercice}[4 points, 15 min]
\begin{enumerate}
    \item En écrivant \( e^{i(\theta + \theta')} \) de deux manières différentes, retrouver les formules d'addition concernant \( \cos(\theta + \theta') \) et \( \sin(\theta + \theta') \).
    \item En déduire que pour tout nombre réel \( \alpha \),
       \[
       \cos^2 \alpha = \dfrac{1 + \cos 2\alpha}{2}
       \]
       et donner une expression de \( \sin 2\alpha \) en fonction de \( \cos \alpha \) et \( \sin \alpha \).
    \item Calculer, à l'aide de la question 2, la valeur de \( \cos\left( \dfrac{\pi}{8} \right) \), puis celle de \( \sin\left( \dfrac{\pi}{8} \right) \).
    %\item Avec une calculatrice, on trouve :
    %   \[
    %   \tan\left( \dfrac{\pi}{12} \right) = 2 - \sqrt{3}.
    %   \]
    %   Démontrer ce résultat à l'aide de la question 2.
\end{enumerate}
\end{exercice}

% Exercice 18
\begin{exercice}[6 points, 25 min]
On dit qu'un triangle équilatéral $ABC$ est direct si et seulement si $(\overrightarrow{AB}; \overrightarrow{AC}) = \frac{\pi}{3}$.
Soient \( A \), \( B \) et \( C \) trois points non alignés d’affixe \( a \), \( b \) et \( c \). On note \( j = e^{\frac{2i\pi}{3}} \).
\begin{enumerate}
    \item \begin{enumerate}
        \item Vérifier que $1, j$ et $j^2$ sont les solutions de l'équation $z^3 = 1$.
        \item Calculer $(1-j)(1 + j + j^2)$. En déduire que $1 + j + j^2 = 0$.
        \item Vérifier que $e^{i\frac{\pi}{3}} + j^2 = 0$.
    \end{enumerate}
    \item \begin{enumerate}
        \item Démontrer que le triangle $ABC$ est équilatéral direct si et seulement si : $\frac{c - a}{b - a} = e^{i\frac{\pi}{3}}$.
        \item Montrer que le triangle \( ABC \) est équilatéral direct si, et seulement si,
        \[
        a + bj + cj^2 = 0.
        \]
    \end{enumerate}
    \item \textbf{(Bonus)} On ne suppose pas nécessairement que \( ABC \) est équilatéral. On construit à partir de \( ABC \) les trois triangles équilatéraux de base \([AB]\), \([AC]\) et \([BC]\) construits à l’extérieur de \( ABC \). Montrer que les centres de gravité de ces trois triangles forment un triangle équilatéral.
\end{enumerate}
\end{exercice}

%% Exercice 19
%\begin{exercice}
%Le plan complexe est muni d'un repère orthonormé \( (O ; \vec{u}, \vec{v}) \) et on considère les points \( A \), \( B \) et \( C \) d'affixes respectives :  
%\( z_A = \sqrt{2}e^{-\frac{3i\pi}{4}} \), \( z_B = 3 + 3i \) et \( z_C = 7 - i \).
%
%\begin{enumerate}
%    \item Déterminer la forme algébrique de \( z_A \) et une forme exponentielle de \( z_B \).
%    \item Quelle est la nature du triangle \( ABC \) ?
%    \item Démontrer qu'il existe un point \( D \) tel que le quadrilatère \( ABCD \) soit un carré. Déterminer l'affixe du point \( D \).
%    \item Déterminer l'affixe du centre du carré \( ABCD \).
%    \item On considère l'ensemble des points \( M \) du plan d'affixe \( z \) tels que \( |z - (3 - i)| = 4 \). Démontrer que cet ensemble est le cercle circonscrit au triangle \( ABC \).
%\end{enumerate}
%\end{exercice}
\end{document}